\section{Related work}\label{sec:related}

\subsection{Event detection on the 3W dataset}

The 3W dataset was introduced with benchmarks for detecting eight types of undesirable events
\citep{vargas2019}, and version 2.0.0 extends it to nine event types and 2,228 instances \citep{vargas2026}.
Detection approaches range from autoencoders combined with rule-based decision diagrams, which report
true-positive rates above 90\,\% in the early stages of an anomaly \citep{aranha2024}, to Transformer
detectors for intelligent completions \citep{aranha2025} and open-world pipelines that combine
reconstruction, binary classifiers and clustering \citep{lopes2025}. A recent layer of large-language-model
agents on top of such a pipeline reports a novelty-detection rate of 94.9\,\% for choke scaling on 3W
segments \citep{lopes2026}. These studies recognise an event that has begun, but none outputs a lead time,
and most do not separate wells between training and test sets; a held-out-well benchmark for 3W has only
recently been proposed \citep{issue231}.

\subsection{Data-driven prediction of mineral scale}

Outside 3W, deep networks detect scale build-up and estimate its magnitude from simulated well histories with
more than 95\,\% detection accuracy \citep{cheng2022}, artificial neural networks predict sulphide-scale
deposition from water analyses and operating data \citep{askar2025}, back-propagation networks estimate the
risk of wellbore blockage from water chemistry in tight gas fields \citep{su2024}, and machine-learning
surrogates reproduce the saturation indices of PHREEQC \citep{parkhurst2013} from a few water properties
\citep{molina2025}. These works establish that scale is predictable from chemistry or from simulated
histories, but they do not use real production-sensor streams, do not forecast the time to a critical
state and are not evaluated on unseen wells.

\subsection{Sequence models and time-to-event learning}

Recurrent networks such as the long short-term memory (LSTM) network \citep{hochreiter1997} and the gated
recurrent unit (GRU) \citep{cho2014}, temporal convolutional networks \citep{bai2018}, attention models
\citep{vaswani2017}, patch-based Transformers \citep{nie2023} and inverted Transformers \citep{liu2024itrans}
are standard tools for multivariate sensor series. A hybrid Transformer--LSTM forecasts water breakthrough in
gas wells two to six days ahead and transfers between similar wells \citep{lai2026}, which makes it the
closest forecasting study of a well event with a lead time. Discrete-time survival heads
\citep{gensheimer2019,kvamme2021} and their deep variants \citep{lee2018deephit} handle censoring and give
monotone risk curves, but they have not been applied to the forecasting of well events. Physics-informed
learning injects domain knowledge through features, architecture or loss terms
\citep{karniadakis2021,willard2022}, while deep ensembles \citep{lakshminarayanan2017} and conformal
prediction \citep{angelopoulos2023} provide uncertainty estimates.

Table~\ref{tab:related} summarises the closest studies. The present work differs from them in four
respects: it forecasts the time to a critical level rather than detecting an event, it uses real
production-sensor data of the choke, it evaluates on wells excluded from training with thresholds fixed in
advance, and it reports calibrated, interval-valued uncertainty.

\begin{table}[t]
\caption{Summary of the closest related studies}
\label{tab:related}
\footnotesize\setlength{\tabcolsep}{3pt}
\begin{tabular*}{\textwidth}{@{\extracolsep\fill}p{2.5cm}p{2.6cm}p{1.9cm}cccp{1.5cm}}
\toprule
Study & Task & Data & Sens. & TTE & HO & Unc. \\
\midrule
\citet{vargas2019,vargas2026} & Event-detection benchmark & 3W & yes & no & no & no \\
\citet{aranha2024} & Anomaly detection & Field wells & yes & no & n/s & no \\
\citet{lopes2025} & Open-world detection & Well completions & yes & no & n/s & no \\
\citet{lopes2026} & Detection, explanation & 3W & yes & no & no & Wilson CI \\
3W proposal \citep{issue231} & Classification benchmark & 3W & yes & no & yes & Bootstrap \\
\citet{cheng2022} & Scale magnitude & Simulated wells & yes & no & no & no \\
\citet{askar2025} & Sulphide-scale risk & Field chemistry & no & no & n/s & no \\
\citet{molina2025} & Saturation-index surrogate & Produced water & no & no & n/s & no \\
\citet{lai2026} & Breakthrough forecast & 9 gas wells & yes & days & partly & no \\
This work & Time to $+5\,\%$ restriction & 3W & yes & yes & yes & Ensemble, conformal \\
\botrule
\end{tabular*}
\footnotetext{Sens.: production-sensor data; TTE: time-to-event forecast; HO: whole wells held out of training; n/s: not stated in the
accessible abstract; Unc.: uncertainty quantification.}
\end{table}
